\chapter{Điểm Số Không Nói Hết Con Người}
\markboth{Điểm Số Không Nói Hết Con Người}{Scores Do Not Describe the Whole Person}

\section{Điểm Thấp Không Nói Hết Nỗ Lực}
\begin{truyen}{Điểm Thấp Không Nói Hết Nỗ Lực}{Low Scores Do Not Erase Effort}
\chuhoa{C}{ó em ôn thật, hỏi thật, làm thật, nhưng bài kiểm tra vẫn không ra như mong muốn.}
\textit{(Some students truly study, ask, and try, yet the test still does not go well.)}

Điểm thấp làm cả lớp và cả nhà dễ quên mất quá trình em đã vật lộn.
\textit{(A low score can make classmates and family forget the struggle that came before it.)}

Một kết quả không đẹp có thể làm mờ cả một quãng cố gắng rất thật.
\textit{(One poor result can blur a very real stretch of effort.)}

Nhưng với người dạy biết nhìn, nỗ lực đó không hề vô nghĩa.
\textit{(But to a teacher who sees carefully, that effort is never meaningless.)}

\end{truyen}

\begin{baihoc}
Điểm số cần được nhìn, nhưng quá trình cũng cần được gọi tên.
\textit{(Scores should be seen, but process should also be named.)}
\end{baihoc}

\begin{ghinhoanh}
\textbf{Short English to remember:}

Scores matter, but process also deserves to be named.

\end{ghinhoanh}

\begin{tuhocnhanh}
\textbf{Gợi ý để dạy và tự nghĩ / Teaching and reflection prompts}

1. Điểm thấp đã che đi điều gì? \textit{What has the low score hidden?}

2. Vì sao quá trình cần được gọi tên? \textit{Why should the process be acknowledged?}

3. Mình đang sửa kết quả hay đang nuôi nỗ lực? \textit{Are we only correcting results or also building effort?}
\end{tuhocnhanh}
\ngancach

\section{Điểm Cao Không Nói Hết Nhân Cách}
\begin{truyen}{Điểm Cao Không Nói Hết Nhân Cách}{High Scores Do Not Guarantee Character}
\chuhoa{C}{ó em học rất giỏi nhưng cư xử với bạn lại lạnh, thiếu tôn trọng, hoặc chỉ biết hơn thua.}
\textit{(Some students learn very well yet treat peers coldly, disrespectfully, or competitively.)}

Nhà trường thường khen người có điểm cao trước.
\textit{(Schools often praise high scorers first.)}

Nhưng nếu chỉ nhìn bảng điểm mà bỏ qua cách làm người, ta rất dễ dạy lệch điều quan trọng.
\textit{(But if we only look at report cards and ignore character, we teach the wrong center.)}

Giỏi chưa bao giờ tự động đồng nghĩa với tử tế.
\textit{(Excellence has never automatically meant kindness.)}

\end{truyen}

\begin{baihoc}
Giáo dục không thể chỉ sản xuất điểm cao. Nó còn phải giữ cho một đứa trẻ biết sống với người khác.
\textit{(Education cannot exist only to produce high scores. It must also preserve a child’s ability to live well with others.)}
\end{baihoc}

\begin{ghinhoanh}
\textbf{Short English to remember:}

Education must preserve a child’s ability to live well with others.

\end{ghinhoanh}

\begin{tuhocnhanh}
\textbf{Gợi ý để dạy và tự nghĩ / Teaching and reflection prompts}

1. Điều gì dễ bị bỏ qua khi chỉ nhìn điểm cao? \textit{What gets overlooked when we focus only on high scores?}

2. Vì sao cần dạy cả cách sống? \textit{Why must education teach character too?}

3. Mình đang khen điều gì nhiều nhất? \textit{What are we praising most often?}
\end{tuhocnhanh}
\ngancach

\section{Một Bài Kém Không Định Nghĩa Một Đứa Trẻ}
\begin{truyen}{Một Bài Kém Không Định Nghĩa Một Đứa Trẻ}{One Bad Test Does Not Define a Child}
\chuhoa{C}{ó hôm một em làm rất tệ. Cả em và người lớn quanh em đều hoảng như thể mọi thứ đã nói hết về con người em.}
\textit{(There are days when a student performs very poorly, and both the child and the adults panic as if the whole person has been defined.)}

Nhưng một bài kém có thể đến từ một ngày xấu, một tuần xấu, hoặc một giai đoạn xấu.
\textit{(Yet one bad test may come from a bad day, a bad week, or a bad season.)}

Trẻ con rất dễ đồng nhất bản thân với kết quả gần nhất.
\textit{(Children easily merge their identity with the latest result.)}

Người lớn cần giúp các em tách hai thứ đó ra.
\textit{(Adults need to help them separate those two things.)}

\end{truyen}

\begin{baihoc}
Kết quả là thông tin, không phải căn cước.
\textit{(A result is information, not identity.)}
\end{baihoc}

\begin{ghinhoanh}
\textbf{Short English to remember:}

A result is information, not identity.

\end{ghinhoanh}

\begin{tuhocnhanh}
\textbf{Gợi ý để dạy và tự nghĩ / Teaching and reflection prompts}

1. Vì sao trẻ dễ tự đồng nhất với điểm số? \textit{Why do children identify themselves with grades so easily?}

2. Người lớn cần nói câu gì vào lúc đó? \textit{What sentence do adults need to say at that moment?}

3. Mình đang sửa một bài hay cứu một lòng tự trọng? \textit{Are we correcting a test or protecting self-worth?}
\end{tuhocnhanh}
\ngancach

\section{Điểm Số Là Tín Hiệu, Không Phải Toàn Bộ}
\begin{truyen}{Điểm Số Là Tín Hiệu, Không Phải Toàn Bộ}{Scores Are Signals, Not the Whole Story}
\chuhoa{N}{hìn đúng, điểm số vẫn có ích. Nó báo cho người dạy biết em đang vướng ở đâu, phần nào chưa chắc, phần nào cần cứu.}
\textit{(Seen correctly, scores are still useful. They signal where a student is stuck, weak, or in need of help.)}

Vấn đề không nằm ở điểm số, mà ở chỗ người lớn dùng nó như dao hay như đèn.
\textit{(The problem is not the score itself, but whether adults use it like a knife or like a lamp.)}

Nếu dùng như dao, nó làm đau. Nếu dùng như đèn, nó soi đường sửa.
\textit{(Used like a knife, it wounds. Used like a lamp, it shows the way to repair.)}

Một con số không đáng sợ bằng cách người lớn đặt ý nghĩa lên nó.
\textit{(A number is less dangerous than the meaning adults attach to it.)}

\end{truyen}

\begin{baihoc}
Điểm số nên được dùng để hiểu học sinh hơn, không phải để làm học sinh nhỏ đi.
\textit{(Scores should be used to understand students better, not to make them feel smaller.)}
\end{baihoc}

\begin{ghinhoanh}
\textbf{Short English to remember:}

Scores should help us understand students better, not make them feel smaller.

\end{ghinhoanh}

\begin{tuhocnhanh}
\textbf{Gợi ý để dạy và tự nghĩ / Teaching and reflection prompts}

1. Mình đang dùng điểm số như dao hay như đèn? \textit{Are we using scores like a knife or a lamp?}

2. Vì sao ý nghĩa gắn lên điểm số quan trọng? \textit{Why does the meaning attached to scores matter so much?}

3. Học sinh cần nghe điều gì sau một bài kiểm tra? \textit{What do students need to hear after a test?}
\end{tuhocnhanh}
\ngancach

\section{Một Điểm Kém Không Xóa Được Một Người Tử Tế}
\begin{truyen}{Một Điểm Kém Không Xóa Được Một Người Tử Tế}{One Bad Grade Cannot Erase a Kind Person}
\chuhoa{C}{ó em học không mạnh, nhưng lúc nào cũng nhặt giùm bạn cây bút rơi, kéo ghế giúp bạn, và nhớ chào thầy cô.}
\textit{(Some students are not strong academically, yet always pick up fallen pens, pull chairs for friends, and greet teachers with care.)}

Nếu chỉ nhìn điểm, người ta dễ bỏ lỡ một con người rất đàng hoàng.
\textit{(If we only look at grades, we can easily miss a very decent human being.)}

Nhà trường cần kết quả, nhưng cũng không thể quên cách một đứa trẻ đối xử với người khác.
\textit{(Schools need results, but cannot forget how a child treats others.)}

Có những giá trị không nằm trong cột điểm nào cả.
\textit{(Some values do not live in any grade column.)}

\end{truyen}

\begin{baihoc}
Điểm số đo được một phần học tập, không đo hết phẩm chất làm người.
\textit{(Grades measure part of learning, not the whole of human character.)}
\end{baihoc}

\begin{ghinhoanh}
\textbf{Short English to remember:}

Grades measure part of learning, not the whole of human character.

\end{ghinhoanh}

\begin{tuhocnhanh}
\textbf{Gợi ý để dạy và tự nghĩ / Teaching and reflection prompts}

1. Em này đang có điều gì mà điểm không đo được? \textit{What does this student have that grades cannot measure?}

2. Vì sao nhà trường dễ bỏ sót phần này? \textit{Why does school easily overlook this part?}

3. Mình có đang ghi nhận điều tốt ngoài điểm không? \textit{Are we recognizing goodness beyond grades?}
\end{tuhocnhanh}
\ngancach

\section{Điểm Cao Nhưng Cách Làm Không Đẹp}
\begin{truyen}{Điểm Cao Nhưng Cách Làm Không Đẹp}{High Scores Achieved in Unhealthy Ways}
\chuhoa{C}{ó em điểm cao thật, nhưng là nhờ học trong hoảng, chép trong sợ, hoặc sống trong áp lực quá mức.}
\textit{(Some students truly score high, yet do it through panic, fear, or excessive pressure.)}

Kết quả đẹp làm người lớn dễ hài lòng.
\textit{(A strong result makes adults feel satisfied quickly.)}

Nhưng nếu cách đi tới kết quả làm một đứa trẻ méo đi, thì điểm đẹp đó cũng đáng để nhìn lại.
\textit{(But if the path to that result distorts a child, then the beautiful score deserves to be reconsidered.)}

Không phải cứ cao là đã khỏe.
\textit{(High does not always mean healthy.)}

\end{truyen}

\begin{baihoc}
Cần nhìn cả kết quả lẫn cái giá mà học sinh đã trả để có kết quả đó.
\textit{(We need to see both the result and the price the student paid for it.)}
\end{baihoc}

\begin{ghinhoanh}
\textbf{Short English to remember:}

We must see both the result and the price paid for it.

\end{ghinhoanh}

\begin{tuhocnhanh}
\textbf{Gợi ý để dạy và tự nghĩ / Teaching and reflection prompts}

1. Điểm cao này có cái giá nào phía sau? \textit{What price may stand behind this high score?}

2. Vì sao “cao” chưa chắc “khỏe”? \textit{Why is “high” not always “healthy”?}

3. Người dạy cần sửa cách học hay chỉ khen kết quả? \textit{Should the teacher adjust the process or only praise the result?}
\end{tuhocnhanh}
\ngancach

\section{Đứa Trẻ Nhìn Điểm Xong Nhìn Mặt Người Lớn}
\begin{truyen}{Đứa Trẻ Nhìn Điểm Xong Nhìn Mặt Người Lớn}{A Child Checks the Grade, Then Checks the Adult’s Face}
\chuhoa{C}{ó những em nhận bài xong không nhìn bạn, không nhìn lại lỗi sai trước. Các em nhìn mặt thầy cô, nhìn mặt cha mẹ.}
\textit{(Some children receive a paper and do not first look at the mistakes or their friends. They look at the adults’ faces.)}

Điều các em tìm không chỉ là kết quả.
\textit{(They are not only checking the result.)}

Các em đang tìm xem với điểm số này, mình còn được thương như cũ không.
\textit{(They are checking whether this score changes how lovable they remain.)}

Điểm số đôi khi đi rất sâu vào nhu cầu được chấp nhận.
\textit{(Grades can reach deeply into the need to be accepted.)}

\end{truyen}

\begin{baihoc}
Sau một bài kiểm tra, phản ứng của người lớn có thể ở lại lâu hơn cả con điểm.
\textit{(After a test, the adult’s reaction may stay longer than the number itself.)}
\end{baihoc}

\begin{ghinhoanh}
\textbf{Short English to remember:}

An adult’s reaction may stay longer than the grade itself.

\end{ghinhoanh}

\begin{tuhocnhanh}
\textbf{Gợi ý để dạy và tự nghĩ / Teaching and reflection prompts}

1. Vì sao trẻ nhìn mặt người lớn trước? \textit{Why do children check the adult’s face first?}

2. Điều gì đang được kiểm tra ngoài con điểm? \textit{What else is being tested besides the number?}

3. Mình muốn học sinh nhớ nét mặt nào của mình? \textit{What expression do we want students to remember?}
\end{tuhocnhanh}
\ngancach

\section{Con Số Cần Đứng Đúng Chỗ}
\begin{truyen}{Con Số Cần Đứng Đúng Chỗ}{The Number Must Stay in Its Proper Place}
\chuhoa{C}{ó giai đoạn người lớn để điểm số ngồi lên quá nhiều thứ: lòng tự trọng, tình thương, hi vọng, và cả cách một đứa trẻ tự gọi tên mình.}
\textit{(There are times adults let grades sit on too many things: self-worth, love, hope, and the way a child names the self.)}

Một con số khi bị đặt quá cao sẽ đè xuống quá nhiều phần khác.
\textit{(A number placed too high can press down too many other parts of life.)}

Điểm số cần thiết. Nhưng nó cần đứng đúng chỗ của nó.
\textit{(Grades matter. But they must stay in their proper place.)}

Khi con số được đặt lại đúng chỗ, việc học mới có cơ hội lành hơn.
\textit{(When the number is put back in place, learning becomes healthier again.)}

\end{truyen}

\begin{baihoc}
Điểm số nên là công cụ của giáo dục, không nên là ông chủ của một đứa trẻ.
\textit{(Grades should be a tool of education, not the master of a child.)}
\end{baihoc}

\begin{ghinhoanh}
\textbf{Short English to remember:}

Grades should be a tool of education, not the master of a child.

\end{ghinhoanh}

\begin{tuhocnhanh}
\textbf{Gợi ý để dạy và tự nghĩ / Teaching and reflection prompts}

1. Con số này đang ngồi lên những điều gì quá mức? \textit{What is this number sitting on too heavily?}

2. Vì sao điểm cần đứng đúng chỗ? \textit{Why must grades stay in their proper place?}

3. Mình có thể đặt lại con số này bằng cách nào? \textit{How can we put the number back in place?}
\end{tuhocnhanh}
